% Essential Formulas

\textbf{Combinatorics \& Modular Queries:}
$C_n^k = \frac{n!}{k!(n-k)!} \pmod M$. Precompute factorials in $O(N)$. Linear inverse:
\lstinline|inv[i] = -(M / i) * inv[M % i] % M + M;|

\textbf{Vandermonde's Identity:} $\sum_{k=0}^r C_m^k C_n^{r-k} = C_{m+n}^r$ (Choosing $r$ items from two combined groups).

\textbf{Fermat's Polygonal Number (Gauss EUREKA):}
Every $N > 0$ can be expressed as the sum of at most $k$ polygonal numbers of side $k$ ($\max \text{count} = k$).
General formula for $n$-th polygonal number of side $k$: 
$$P_k(n) = \frac{(k - 2)n^2 - (k - 4)n}{2} = n + (k - 2) \cdot \frac{n(n-1)}{2}$$
\begin{itemize}
    \item \textbf{Recurrence Relation between Sides:} A $k$-polygonal number can be derived from a $(k-1)$-polygonal number by adding the $(n-1)$-th triangular number ($T_{n-1} = \frac{n(n-1)}{2}$):
    $$P_k(n) = P_{k-1}(n) + T_{n-1} = P_{k-1}(n) + \frac{n(n-1)}{2}$$
    \item \textbf{Triangular Decomposition:} Any $k$-polygonal number can be broken down directly into a combination of the $n$-th and $(n-1)$-th triangular numbers:
    $$P_k(n) = T_n + (k - 3) \cdot T_{n-1}$$
    \item \textbf{General $O(1)$ Check / Inverse Function:} To check if $N$ is a $k$-polygonal number, solve $(k - 2)n^2 - (k - 4)n - 2N = 0$. Let $\Delta = (k - 4)^2 + 8N(k - 2)$. $N$ is valid iff $\Delta$ is a perfect square and $n = \frac{(k - 4) + \sqrt{\Delta}}{2(k - 2)}$ is an integer.
    \item \textbf{Find side $k$ from $N$ and $n$:} $k = 2 + \frac{2(N - n)}{n(n - 1)}$.
\end{itemize}

\textbf{Catalan Numbers:} (BSTs, valid parenthesization, non-crossing polygon dissections).
$C_0 = 1, \, C_n = \sum_{i=0}^{n-1} C_i C_{n-1-i} = \frac{1}{n+1} C_{2n}^n = C_{2n}^n - C_{2n}^{n-1}$
Terms: $1, 1, 2, 5, 14, 42, 132, 429, 1430, 4862$.

\textbf{Derangements:} (Permutations with no fixed points).
$D_1 = 0, \, D_2 = 1, \, D_n = (n-1)(D_{n-1} + D_{n-2}) = n D_{n-1} + (-1)^n$
Terms: $0, 1, 2, 9, 44, 265, 1854, 14833, 133496$.

\textbf{Stirling Numbers of the 2nd Kind:} (Partition $n$ labeled items into $k$ unlabeled non-empty subsets).
$S(n, k) = S(n-1, k-1) + k \cdot S(n-1, k)$. Base: $S(0, 0) = 1, \, S(n, 0) = S(0, n) = 0$.
Partition $n$ labeled items into $k$ labeled boxes $= k! \cdot S(n, k)$.

\textbf{Lucas' Theorem:} $C_n^k \pmod p$ for small prime $p$ ($n, k \le 10^{18}$). Base-$p$ representations:
$n = \sum n_i p^i, \, k = \sum k_i p^i \implies C_n^k \equiv \prod C_{n_i}^{k_i} \pmod p$ (where $C_{n_i}^{k_i} = 0$ if $n_i < k_i$).

\textbf{Burnside's Lemma:} Number of orbits (distinct objects under symmetry) $= \frac{1}{|G|} \sum_{g \in G} |X^g|$ where $|G|$ is size of symmetry group, and $|X^g|$ is number of elements fixed by symmetry $g$.

\textbf{Graph Theory Counting:}
\begin{itemize}
    \item \textbf{Cayley's Formula:} Labeled spanning trees in $K_n = n^{n-2}$. In a forest of $k$ trees with roots $r_1..r_k$ covering all $n$ nodes, ways to connect them into a single tree $= n^{k-2} \cdot \prod \text{sz}(T_i)$.
    \item \textbf{Euler's Planar Formula:} Connected planar graph: $V - E + F = 2$.
    \item \textbf{Matrix Tree Theorem:} Let $A$ be adjacency matrix, $D$ degree matrix. Kirchhoff matrix $L = D - A$. Delete row $i$ and col $i$. Determinant of resulting matrix $=$ number of spanning trees.
\end{itemize}

\textbf{Geometry Theorems:}
\begin{itemize}
    \item \textbf{Pick's Theorem:} Grid polygon area $A = I + B/2 - 1$. ($I$: interior points, $B$: boundary points).
    \item \textbf{Boundary Points:} Integer points on segment $(x_1, y_1)$ to $(x_2, y_2)$ inclusive $=\gcd(|x_1 - x_2|, |y_1 - y_2|) + 1$.
\end{itemize}
